\documentclass[10pt,a4paper]{amsart}
\usepackage{amsmath,amssymb,amsthm,graphicx}
\usepackage{verbatim}
\newenvironment{reference}{\comment}{\endcomment}
%\newenvironment{reference}{}{}
\newenvironment{slogan}{\comment}{\endcomment}
\newenvironment{history}{\comment}{\endcomment}
% For commutative diagrams we use Xy-pic
\usepackage[all]{xy}

% We use 2cell for 2-commutative diagrams.
\xyoption{2cell}
% Theorem environments.
%
\theoremstyle{plain}
\newtheorem{theorem}[subsection]{Theorem}
\newtheorem{proposition}[subsection]{Proposition}
\newtheorem{lemma}[subsection]{Lemma}

\theoremstyle{definition}
\newtheorem{definition}[subsection]{Definition}
\newtheorem{example}[subsection]{Example}
\newtheorem{exercise}[subsection]{Exercise}
\newtheorem{situation}[subsection]{Situation}

\theoremstyle{remark}
\newtheorem{remark}[subsection]{Remark}
\newtheorem{remarks}[subsection]{Remarks}

\numberwithin{equation}{subsection}

% Macros
%
\def\lim{\mathop{\mathrm{lim}}\nolimits}
\def\colim{\mathop{\mathrm{colim}}\nolimits}
\def\Spec{\mathop{\mathrm{Spec}}}
\def\Hom{\mathop{\mathrm{Hom}}\nolimits}
\def\Ext{\mathop{\mathrm{Ext}}\nolimits}
\def\SheafHom{\mathop{\mathcal{H}\!\mathit{om}}\nolimits}
\def\SheafExt{\mathop{\mathcal{E}\!\mathit{xt}}\nolimits}
\def\Sch{\mathit{Sch}}
\def\Mor{\mathop{\mathrm{Mor}}\nolimits}
\def\Ob{\mathop{\mathrm{Ob}}\nolimits}
\def\Sh{\mathop{\mathit{Sh}}\nolimits}
\def\NL{\mathop{N\!L}\nolimits}
\def\CH{\mathop{\mathrm{CH}}\nolimits}
\def\proetale{{pro\text{-}\acute{e}tale}}
\def\etale{{\acute{e}tale}}
\def\QCoh{\mathit{QCoh}}
\def\Ker{\mathop{\mathrm{Ker}}}
\def\Im{\mathop{\mathrm{Im}}}
\def\Coker{\mathop{\mathrm{Coker}}}
\def\Coim{\mathop{\mathrm{Coim}}}

% Boxtimes
%
\DeclareMathSymbol{\boxtimes}{\mathbin}{AMSa}{"02}

%
% Macros for moduli stacks/spaces
%
\def\QCohstack{\mathcal{QC}\!\mathit{oh}}
\def\Cohstack{\mathcal{C}\!\mathit{oh}}
\def\Spacesstack{\mathcal{S}\!\mathit{paces}}
\def\Quotfunctor{\mathrm{Quot}}
\def\Hilbfunctor{\mathrm{Hilb}}
\def\Curvesstack{\mathcal{C}\!\mathit{urves}}
\def\Polarizedstack{\mathcal{P}\!\mathit{olarized}}
\def\Complexesstack{\mathcal{C}\!\mathit{omplexes}}
% \Pic is the operator that assigns to X its picard group, usage \Pic(X)
% \Picardstack_{X/B} denotes the Picard stack of X over B
% \Picardfunctor_{X/B} denotes the Picard functor of X over B
\def\Pic{\mathop{\mathrm{Pic}}\nolimits}
\def\Picardstack{\mathcal{P}\!\mathit{ic}}
\def\Picardfunctor{\mathrm{Pic}}
\def\Deformationcategory{\mathcal{D}\!\mathit{ef}}
\usepackage{cleveref}
\crefname{originallemma}{original lemma}{original lemmas}
\Crefname{originallemma}{Original lemma}{Original lemmas}
\makeatletter
\newlabel{coherent-lemma-quasi-compact-h1-zero-invertible}{{0B5P}{source}{}{sourcepointer-1}{}}
\newlabel{coherent-lemma-quasi-compact-h1-zero-invertible@cref}{{[originallemma][0][]0B5P}{[1][0][]source}{}{}{}}
\makeatother
\crefname{originallemma}{original lemma}{original lemmas}
\Crefname{originallemma}{Original lemma}{Original lemmas}
\makeatletter
\newlabel{decent-spaces-lemma-decent-space-closed-point}{{0AHB}{source}{}{sourcepointer-2}{}}
\newlabel{decent-spaces-lemma-decent-space-closed-point@cref}{{[originallemma][0][]0AHB}{[1][0][]source}{}{}{}}
\makeatother
\crefname{originalsection}{original section}{original sections}
\Crefname{originalsection}{Original section}{Original sections}
\makeatletter
\newlabel{decent-spaces-section-reasonable-decent}{{03I7}{source}{}{sourcepointer-3}{}}
\newlabel{decent-spaces-section-reasonable-decent@cref}{{[originalsection][0][]03I7}{[1][0][]source}{}{}{}}
\makeatother
\crefname{originallemma}{original lemma}{original lemmas}
\Crefname{originallemma}{Original lemma}{Original lemmas}
\makeatletter
\newlabel{lemma-affine-vanishing-higher-direct-images}{{073H}{source}{}{sourcepointer-4}{}}
\newlabel{lemma-affine-vanishing-higher-direct-images@cref}{{[originallemma][0][]073H}{[1][0][]source}{}{}{}}
\makeatother
\crefname{originallemma}{original lemma}{original lemmas}
\Crefname{originallemma}{Original lemma}{Original lemmas}
\makeatletter
\newlabel{lemma-colimits}{{073E}{source}{}{sourcepointer-5}{}}
\newlabel{lemma-colimits@cref}{{[originallemma][0][]073E}{[1][0][]source}{}{}{}}
\makeatother
\crefname{originallemma}{original lemma}{original lemmas}
\Crefname{originallemma}{Original lemma}{Original lemmas}
\makeatletter
\newlabel{lemma-directed-colimit-coherent}{{07UV}{source}{}{sourcepointer-6}{}}
\newlabel{lemma-directed-colimit-coherent@cref}{{[originallemma][0][]07UV}{[1][0][]source}{}{}{}}
\makeatother
\crefname{originallemma}{original lemma}{original lemmas}
\Crefname{originallemma}{Original lemma}{Original lemmas}
\makeatletter
\newlabel{morphisms-lemma-affine-s-open}{{01SF}{source}{}{sourcepointer-7}{}}
\newlabel{morphisms-lemma-affine-s-open@cref}{{[originallemma][0][]01SF}{[1][0][]source}{}{}{}}
\makeatother
\crefname{originalproposition}{original proposition}{original propositions}
\Crefname{originalproposition}{Original proposition}{Original propositions}
\makeatletter
\newlabel{proposition-vanishing-affine}{{07V6}{source}{}{sourcepointer-8}{}}
\newlabel{proposition-vanishing-affine@cref}{{[originalproposition][0][]07V6}{[1][0][]source}{}{}{}}
\makeatother
\crefname{originallemma}{original lemma}{original lemmas}
\Crefname{originallemma}{Original lemma}{Original lemmas}
\makeatletter
\newlabel{sites-cohomology-lemma-apply-Leray}{{0733}{source}{}{sourcepointer-9}{}}
\newlabel{sites-cohomology-lemma-apply-Leray@cref}{{[originallemma][0][]0733}{[1][0][]source}{}{}{}}
\makeatother
\crefname{originallemma}{original lemma}{original lemmas}
\Crefname{originallemma}{Original lemma}{Original lemmas}
\makeatletter
\newlabel{spaces-morphisms-lemma-affine-local}{{03WG}{source}{}{sourcepointer-10}{}}
\newlabel{spaces-morphisms-lemma-affine-local@cref}{{[originallemma][0][]03WG}{[1][0][]source}{}{}{}}
\makeatother
\crefname{originallemma}{original lemma}{original lemmas}
\Crefname{originallemma}{Original lemma}{Original lemmas}
\makeatletter
\newlabel{spaces-properties-lemma-Noetherian-sober}{{04ZG}{source}{}{sourcepointer-11}{}}
\newlabel{spaces-properties-lemma-Noetherian-sober@cref}{{[originallemma][0][]04ZG}{[1][0][]source}{}{}{}}
\makeatother
\crefname{originallemma}{original lemma}{original lemmas}
\Crefname{originallemma}{Original lemma}{Original lemmas}
\makeatletter
\newlabel{spaces-properties-lemma-subscheme}{{03JH}{source}{}{sourcepointer-12}{}}
\newlabel{spaces-properties-lemma-subscheme@cref}{{[originallemma][0][]03JH}{[1][0][]source}{}{}{}}
\makeatother
\title{A cohomological vanishing hypothesis on a Noetherian algebraic space forces it to be a scheme and the invertible sheaf to be ample}
\author{The Stacks Project Authors}
\date{stacks-project:spaces-cohomology.tex:local-snapshot-20261012; GNU FDL 1.2 or later}
\begin{document}
\maketitle
\noindent\textit{Editorial note.} All delivered mathematical text is the author's own, copied from the pinned Stacks Project chapter file spaces-cohomology.tex; no statement or proof step was rewritten, reordered or removed.\par\medskip
\noindent\textit{Editorial note.} Preamble provenance: the pinned chapter file opens with \textbackslash{}input\{preamble\} and carries no preamble text of its own, so the project's single Stacks-wide preamble staged at sources/stacks-spaces-cohomology/preamble.tex is excerpted verbatim as three spans, lines 8-12 (verbatim package and the project's own reference/slogan/history display environments), lines 14-18 (Xy-pic setup) and lines 154-217 (theorem declarations and mathematical macros). The wrapper class and the amsmath/amssymb/amsthm/graphicx base packages are supplied by the assembler. No macro was invented or redefined.\par\medskip
\noindent\textit{Editorial note.} Same-chapter pointers are located at their \textbackslash{}label lines inside the pinned chapter file; cross-chapter pointers are located in the project's own tag list staged in the snapshot (sources/stacks-chow/tags.txt), at the verbatim lines carrying their canonical tags.\par\medskip
\noindent\textit{Allowed standard prerequisites (editorial).}\begin{itemize}
\item A space is affine as soon as the first cohomology of every quasi-coherent module vanishes. Same chapter spaces-cohomology.tex, label proposition-vanishing-affine (invoked at proof line 3248). It reduces affineness of the open \$U\$ below a section to the vanishing \$H\textasciicircum{}1(U, \textbackslash{}mathcal\{F\}) = 0\$ for all quasi-coherent \$\textbackslash{}mathcal\{F\}\$.
\item The inclusion of the nonvanishing open of a section is affine. Morphisms, label morphisms-lemma-affine-s-open (invoked at proof line 3251) and Morphisms of Spaces, label spaces-morphisms-lemma-affine-local (invoked at proof line 3252). They show that \$j : U \textbackslash{}to X\$ is affine, which is what turns \$H\textasciicircum{}1(U, -)\$ into \$H\textasciicircum{}1(X, j\_*-)\$.
\item Vanishing of the higher direct images of an affine morphism, and the Leray identification. Same chapter spaces-cohomology.tex, label lemma-affine-vanishing-higher-direct-images (invoked at proof line 3258) and Cohomology on Sites, label sites-cohomology-lemma-apply-Leray (invoked at proof line 3259). Together they give \$H\textasciicircum{}1(U, \textbackslash{}mathcal\{F\}) = H\textasciicircum{}1(X, j\_*\textbackslash{}mathcal\{F\})\$.
\item A quasi-coherent module as a filtered colimit of coherent modules, and cohomology of quasi-compact quasi-separated spaces commuting with such colimits. Same chapter spaces-cohomology.tex, labels lemma-directed-colimit-coherent (invoked at proof line 3262) and lemma-colimits (invoked at proof line 3266). They replace \$j\_*\textbackslash{}mathcal\{F\}\$ by the filtered colimit of coherent \$\textbackslash{}mathcal\{F\}\_i\$ and compute \$H\textasciicircum{}1(X, j\_*\textbackslash{}mathcal\{F\})\$ as the colimit of the \$H\textasciicircum{}1(X, \textbackslash{}mathcal\{F\}\_i)\$.
\item Closed points of a decent Noetherian space, the schematic locus, and Noetherian sobriety. Decent Spaces, label decent-spaces-lemma-decent-space-closed-point (invoked at proof line 3300), Properties of Spaces, label spaces-properties-lemma-subscheme (invoked at proof line 3320) and spaces-properties-lemma-Noetherian-sober (invoked at proof line 3328). They let every closed point be cut out by a closed immersion from a field, show that the closed points lie in the schematic locus once the \$U\$ are affine, and conclude that an open subset containing all closed points is everything.
\item The scheme case of the criterion. Cohomology of Schemes, label coherent-lemma-quasi-compact-h1-zero-invertible (invoked at proof line 3331). Once \$X\$ is known to be a scheme, it converts the assumed vanishing into ampleness of \$\textbackslash{}mathcal\{L\}\$.
\end{itemize}\medskip
\ifdefined\hypertarget\hypertarget{sourcepointer-1}{}\fi
\noindent\textit{Source reference.} Original lemma 0B5P: sources/stacks-chow/tags.txt, source lines 14042--14042. Navigation pointer for \textbackslash{}ref\{coherent-lemma-quasi-compact-h1-zero-invertible\} used inside the counted span (line 3331); the argument that consumes it is retained in the span.\par\medskip
\ifdefined\hypertarget\hypertarget{sourcepointer-2}{}\fi
\noindent\textit{Source reference.} Original lemma 0AHB: sources/stacks-chow/tags.txt, source lines 13224--13224. Navigation pointer for \textbackslash{}ref\{decent-spaces-lemma-decent-space-closed-point\} used inside the counted span (line 3300); the argument that consumes it is retained in the span.\par\medskip
\ifdefined\hypertarget\hypertarget{sourcepointer-3}{}\fi
\noindent\textit{Source reference.} Original section 03I7: sources/stacks-chow/tags.txt, source lines 4593--4593. Navigation pointer for \textbackslash{}ref\{decent-spaces-section-reasonable-decent\} used inside the counted span (line 3304); the argument that consumes it is retained in the span.\par\medskip
\ifdefined\hypertarget\hypertarget{sourcepointer-4}{}\fi
\noindent\textit{Source reference.} Original lemma 073H: sources/stacks-spaces-cohomology/pinned-original-source.tex, source lines 1522--1522. Navigation pointer for \textbackslash{}ref\{lemma-affine-vanishing-higher-direct-images\} used inside the counted span (line 3258); the argument that consumes it is retained in the span.\par\medskip
\ifdefined\hypertarget\hypertarget{sourcepointer-5}{}\fi
\noindent\textit{Source reference.} Original lemma 073E: sources/stacks-spaces-cohomology/pinned-original-source.tex, source lines 667--667. Navigation pointer for \textbackslash{}ref\{lemma-colimits\} used inside the counted span (line 3266); the argument that consumes it is retained in the span.\par\medskip
\ifdefined\hypertarget\hypertarget{sourcepointer-6}{}\fi
\noindent\textit{Source reference.} Original lemma 07UV: sources/stacks-spaces-cohomology/pinned-original-source.tex, source lines 2866--2866. Navigation pointer for \textbackslash{}ref\{lemma-directed-colimit-coherent\} used inside the counted span (line 3262); the argument that consumes it is retained in the span.\par\medskip
\ifdefined\hypertarget\hypertarget{sourcepointer-7}{}\fi
\noindent\textit{Source reference.} Original lemma 01SF: sources/stacks-chow/tags.txt, source lines 2336--2336. Navigation pointer for \textbackslash{}ref\{morphisms-lemma-affine-s-open\} used inside the counted span (line 3251); the argument that consumes it is retained in the span.\par\medskip
\ifdefined\hypertarget\hypertarget{sourcepointer-8}{}\fi
\noindent\textit{Source reference.} Original proposition 07V6: sources/stacks-spaces-cohomology/pinned-original-source.tex, source lines 3160--3160. Navigation pointer for \textbackslash{}ref\{proposition-vanishing-affine\} used inside the counted span (line 3248); the argument that consumes it is retained in the span.\par\medskip
\ifdefined\hypertarget\hypertarget{sourcepointer-9}{}\fi
\noindent\textit{Source reference.} Original lemma 0733: sources/stacks-chow/tags.txt, source lines 9047--9047. Navigation pointer for \textbackslash{}ref\{sites-cohomology-lemma-apply-Leray\} used inside the counted span (line 3259); the argument that consumes it is retained in the span.\par\medskip
\ifdefined\hypertarget\hypertarget{sourcepointer-10}{}\fi
\noindent\textit{Source reference.} Original lemma 03WG: sources/stacks-chow/tags.txt, source lines 5111--5111. Navigation pointer for \textbackslash{}ref\{spaces-morphisms-lemma-affine-local\} used inside the counted span (line 3253); the argument that consumes it is retained in the span.\par\medskip
\ifdefined\hypertarget\hypertarget{sourcepointer-11}{}\fi
\noindent\textit{Source reference.} Original lemma 04ZG: sources/stacks-chow/tags.txt, source lines 6464--6464. Navigation pointer for \textbackslash{}ref\{spaces-properties-lemma-Noetherian-sober\} used inside the counted span (line 3328); the argument that consumes it is retained in the span.\par\medskip
\ifdefined\hypertarget\hypertarget{sourcepointer-12}{}\fi
\noindent\textit{Source reference.} Original lemma 03JH: sources/stacks-chow/tags.txt, source lines 4639--4639. Navigation pointer for \textbackslash{}ref\{spaces-properties-lemma-subscheme\} used inside the counted span (line 3320); the argument that consumes it is retained in the span.\par\medskip
\begin{lemma}
\label{lemma-Noetherian-h1-zero-invertible}
Let $S$ be a scheme. Let $X$ be a Noetherian algebraic space over $S$.
Let $\mathcal{L}$ be an invertible $\mathcal{O}_X$-module.
Assume that for every coherent $\mathcal{O}_X$-module
$\mathcal{F}$ there exists an $n \geq 1$ such that
$H^1(X, \mathcal{F} \otimes_{\mathcal{O}_X} \mathcal{L}^{\otimes n}) = 0$.
Then $X$ is a scheme and $\mathcal{L}$ is ample on $X$.
\end{lemma}

\begin{proof}
Let $s \in H^0(X, \mathcal{L}^{\otimes d})$ be a global section.
Let $U \subset X$ be the open subspace over which $s$ is a generator
of $\mathcal{L}^{\otimes d}$. In particular we have
$\mathcal{L}^{\otimes d}|_U \cong \mathcal{O}_U$.
We claim that $U$ is affine.

\medskip\noindent
Proof of the claim. We will show that $H^1(U, \mathcal{F}) = 0$
for every quasi-coherent $\mathcal{O}_U$-module $\mathcal{F}$.
This will prove the claim by Proposition \ref{proposition-vanishing-affine}.
Denote $j : U \to X$ the inclusion morphism.
Since \'etale locally the morphism $j$ is affine
(by Morphisms, Lemma \ref{morphisms-lemma-affine-s-open})
we see that $j$ is affine (Morphisms of Spaces, Lemma
\ref{spaces-morphisms-lemma-affine-local}).
Hence we have
$$
H^1(U, \mathcal{F}) = H^1(X, j_*\mathcal{F})
$$
by Lemma \ref{lemma-affine-vanishing-higher-direct-images}
(and Cohomology on Sites, Lemma \ref{sites-cohomology-lemma-apply-Leray}).
Write $j_*\mathcal{F} = \colim \mathcal{F}_i$ as a filtered colimit
of coherent $\mathcal{O}_X$-modules, see
Lemma \ref{lemma-directed-colimit-coherent}. Then
$$
H^1(X, j_*\mathcal{F}) = \colim H^1(X, \mathcal{F}_i)
$$
by Lemma \ref{lemma-colimits}.
Thus it suffices to show that $H^1(X, \mathcal{F}_i)$ maps
to zero in $H^1(U, j^*\mathcal{F}_i)$. By assumption there exists
an $n \geq 1$ such that
$$
H^1(X,
\mathcal{F}_i \otimes_{\mathcal{O}_X}
(\mathcal{O}_X \oplus \mathcal{L} \oplus \ldots
\oplus \mathcal{L}^{\otimes d - 1})
\otimes_{\mathcal{O}_X} \mathcal{L}^{\otimes n}) = 0
$$
Hence there exists an $a \geq 0$ such that
$H^1(X, \mathcal{F}_i \otimes_{\mathcal{O}_X} \mathcal{L}^{\otimes ad}) = 0$.
On the other hand, the map
$$
s^a : \mathcal{F}_i \longrightarrow
\mathcal{F}_i \otimes_{\mathcal{O}_X} \mathcal{L}^{\otimes ad}
$$
is an isomorphism after restriction to $U$. Contemplating the
commutative diagram
$$
\xymatrix{
H^1(X, \mathcal{F}_i) \ar[r] \ar[d]_{s^a} & H^1(U, j^*\mathcal{F}_i)
\ar[d]^{\cong} \\
H^1(X, \mathcal{F}_i \otimes_{\mathcal{O}_X} \mathcal{L}^{\otimes ad}) \ar[r] &
H^1(U,
j^*(\mathcal{F}_i \otimes_{\mathcal{O}_X} \mathcal{L}^{\otimes ad}))
}
$$
we conclude that the map  $H^1(X, \mathcal{F}_i) \to H^1(U, j^*\mathcal{F}_i)$
is zero and the claim holds.

\medskip\noindent
Let $x \in |X|$ be a closed point. By Decent Spaces, Lemma
\ref{decent-spaces-lemma-decent-space-closed-point}
we can represent $x$ by a closed immersion $i : \Spec(k) \to X$
(this also uses that a quasi-separated algebraic space is
decent, see Decent Spaces, Section
\ref{decent-spaces-section-reasonable-decent}).
Thus $\mathcal{O}_X \to i_*\mathcal{O}_{\Spec(k)}$ is surjective.
Let $\mathcal{I} \subset \mathcal{O}_X$ be the kernel and choose
$d \geq 1$ such that
$H^1(X, \mathcal{I} \otimes_{\mathcal{O}_X} \mathcal{L}^{\otimes d}) = 0$.
Then
$$
H^0(X, \mathcal{L}^{\otimes d}) \to
H^0(X,
i_*\mathcal{O}_{\Spec(k)} \otimes_{\mathcal{O}_X} \mathcal{L}^{\otimes d}) =
H^0(\Spec(k), i^*\mathcal{L}^{\otimes d}) \cong k
$$
is surjective by the long exact cohomology sequence. Hence
there exists an $s \in H^0(X, \mathcal{L}^{\otimes d})$
such that $x \in U$ where $U$ is the open subspace corresponding to $s$
as above. Thus $x$ is in the schematic locus
(see Properties of Spaces, Lemma \ref{spaces-properties-lemma-subscheme})
of $X$ by our claim.

\medskip\noindent
To conclude that $X$ is a scheme, it suffices to show that
any open subset of $|X|$ which contains all the closed points
is equal to $|X|$. This follows from the fact that $|X|$
is a Noetherian topological space, see
Properties of Spaces, Lemma \ref{spaces-properties-lemma-Noetherian-sober}.
Finally, if $X$ is a scheme, then we can apply
Cohomology of Schemes, Lemma
\ref{coherent-lemma-quasi-compact-h1-zero-invertible}
to conclude that $\mathcal{L}$ is ample.
\end{proof}

\end{document}
